% !TeX program = lualatex
% Compile twice: lualatex 114.tex
% All references and prose are embedded; no BibTeX database or external inputs.
\documentclass[11pt,a4paper]{article}
\usepackage[margin=24mm,headheight=14pt]{geometry}
\usepackage{fontspec}
\setmainfont{Latin Modern Roman}
\setsansfont{Latin Modern Sans}
\setmonofont{Latin Modern Mono}
\usepackage{amsmath,amssymb,mathtools}
\usepackage{unicode-math}
\setmathfont{Latin Modern Math}
\usepackage{microtype}
\usepackage{enumitem,xurl,needspace}
\usepackage[dvipsnames]{xcolor}
\usepackage{hyperref}
\hypersetup{unicode=true,colorlinks=true,linkcolor=MidnightBlue,urlcolor=MidnightBlue,
 pdftitle={Research attempt: Problem 114},pdfauthor={Research notebook},bookmarksnumbered=true}
\usepackage{bookmark}
\usepackage{tocloft}
\setlength{\cftsecnumwidth}{3em}
\cftsetpnumwidth{2.5em}
\cftsetrmarg{4em}
\usepackage{titlesec}
\titleformat{\section}{\Large\bfseries\raggedright}{\thesection}{0.7em}{}
\usepackage{fancyhdr}
\pagestyle{fancy}\fancyhf{}
\fancyhead[L]{\small\sffamily Open Problems Math | Research notebook}
\fancyhead[R]{\small Problem 114 | 27 September 2026}
\fancyfoot[C]{\thepage}
\setlength{\parindent}{0pt}
\setlength{\parskip}{3pt plus 1pt}
\setlength{\emergencystretch}{3em}
\setcounter{tocdepth}{1}
\setcounter{secnumdepth}{1}
\setlist[itemize]{leftmargin=3em,itemsep=5pt,topsep=4pt,parsep=0pt}
\allowdisplaybreaks
\newcommand{\N}{\mathbb{N}}
\newcommand{\Z}{\mathbb{Z}}
\newcommand{\Q}{\mathbb{Q}}
\newcommand{\R}{\mathbb{R}}
\newcommand{\C}{\mathbb{C}}
\newcommand{\F}{\mathbb{F}}
\newcommand{\A}{\mathbb{A}}
\newcommand{\PP}{\mathbb P}
\newcommand{\E}{\mathbb{E}}
\newcommand{\eps}{\varepsilon}
\newcommand{\dd}{\,\mathrm d}
\newcommand{\sref}[2]{\hyperlink{source:#1:#2}{[S#2]}}
\newcommand{\eref}[2]{\hyperlink{extra:#1:#2}{[E#2]}}
% Turn the next switch off for a shorter reading copy. The quotations preserve
% every exactTarget field, including qualifications omitted from a short summary.
\newif\ifshowcatalogue
\showcataloguetrue
\newcommand{\cataloguescope}[1]{\ifshowcatalogue
 \paragraph{Exact scope in the supplied catalogue.}
 \begin{quote}\small #1\end{quote}\fi}
\usepackage{etoolbox}
\AtBeginEnvironment{itemize}{\small}
\numberwithin{equation}{section}
% Standalone export. Original section 114 preserved verbatim outside marked additions.
\begin{document}
\setcounter{section}{113}
% ================================================================
% problemId: problem.iitaka-conjecture-c-n-m-on-subadditivity-of-kodaira-dimension
% source releaseRank: 114; source releaseStatus: open_with_solved_subcases
% exactTarget SHA-256: 6fd32a503735df46a7865bf616f52174f7303db0fc51fa8e4da065c031e53ba1
\clearpage
\section{\texorpdfstring{Iitaka conjecture C\_\{n,m\} on subadditivity of Kodaira dimension}{Iitaka conjecture C\_\{n,m\} on subadditivity of Kodaira dimension}}\label{114}
\subsection{Definitions and mathematical statement}
For a smooth projective variety $Z$, its Kodaira dimension $\kappa(Z)$ is the growth exponent of the spaces $H^0(Z,mK_Z)$, with value $-\infty$ if all positive pluricanonical spaces vanish. For the surjective morphism and general fiber in the record, the target is
\[
 \kappa(X)\ge\kappa(Y)+\kappa(F).
\]
The usual fibration convention uses connected fibers and the characteristic-zero setting of the cited source; the catalogue's abbreviated sentence does not separately spell out those conventions. When one term on the right is $-\infty$, the inequality is interpreted in the extended ordered sense.
\par\smallskip\textit{Formulation and concept references:} \sref{114}{1}.
\cataloguescope{Prove the Iitaka conjecture C\_\{n,m\}: for a surjective morphism f: X -> Y between smooth projective varieties with general fibre F, kappa(X) >= kappa(Y) + kappa(F).}
\subsection{Short English statement}
Is the Kodaira dimension of a projective fibration at least the sum of those of its base and general fiber?
% BEGIN RESEARCH ADDITION 114
\subsection{Research attempt}

\paragraph{Current frontier and source audit.}
Birkar's introduction explicitly works over an algebraically closed field
of characteristic zero and defines an algebraic fibre space to have
connected fibres. Those are the conventions incorporated by the present
entry. Theorem~1.3 proves the conjecture for $\dim X\leq6$; the introduction
also records the good-minimal-model fibre case. Neither unrestricted
existence of good minimal models nor abundance is assumed here.
\sref{114}{1}

Established methods include positivity of pluricanonical direct images;
Cao--P\u{a}un prove the abelian-base case, with a logarithmic extension.
A December 2025 manuscript by Ammar states the case
$\dim\operatorname{alb}_Y(Y)\geq\dim Y-2$. That recent result is reported
as a manuscript theorem, not used as an independently audited ingredient.
The complete arbitrary-base, arbitrary-fibre claim is not supplied by
these hypotheses. \eref{114}{1}\eref{114}{2}

\paragraph{Attack route.}
A basis of pluricanonical forms on the generic fibre extends meromorphically
over the base. The central question is the cost of removing its poles.
The route below turns this into a precise section-counting lemma, then
asks whether positivity can bound that cost in directions supported by
the base's canonical system. A second route tests descent through a
finite \'etale product cover. A third, numerical-positivity-only route is
rejected by a degree-zero line-bundle example.

\paragraph{Attempt: a pole-cost criterion.}
Write $K_{X/Y}=K_X-f^*K_Y$ and
\[
 E_m=f_*\mathcal O_X(mK_{X/Y}),\qquad
 r_m=\operatorname{rank}E_m=h^0(F,mK_F).
\]
The last equality is generic base change; the general fibre is understood
in that sense for each $m$. The sheaf is torsion-free, and all generic-rank
arguments below take place on a dense open subset before being interpreted
as morphisms of coherent sheaves on $Y$. Projection formula gives
\begin{equation}\label{eq:114:projection}
 h^0(X,mK_X)=h^0(Y,E_m\otimes\mathcal O_Y(mK_Y)).
\end{equation}
If either Kodaira dimension on the right of the target is $-\infty$, the
inequality is automatic. Assume from now on that both are nonnegative.

\textbf{Lemma 1 (canonical-direction pole cost).}
Choose an integer $q>0$ and an effective divisor $D\sim qK_Y$.
Suppose along an unbounded sequence of multiples $m$ of $q$ there are
integers $b_m\geq0$ with $b_m=o(m)$ and morphisms
\begin{equation}\label{eq:114:injection}
 \mathcal O_Y(-b_mD)^{\oplus r_m}\hookrightarrow E_m
\end{equation}
that have rank $r_m$ at the generic point. Choose $q$ sufficiently divisible
that the standard lower plurigenus-growth bounds for $F$ and $Y$ hold on
its multiples. Then $\kappa(X)\geq\kappa(F)+\kappa(Y)$.

\emph{Proof.} A generically injective map from the displayed locally free
sheaf is injective everywhere as a sheaf map: its kernel is torsion-free
and zero at the generic point. Tensoring by $\mathcal O_Y(mK_Y)$ and taking
sections gives
\[
 h^0(X,mK_X)\geq r_m\,
 h^0\bigl(Y,(m-qb_m)K_Y\bigr).
\]
Eventually $m-qb_m\geq m/2$ and it is still a multiple of $q$. Thus the
right side is at least $c m^{\kappa(F)+\kappa(Y)}$ along the sequence.
This yields the asserted Kodaira dimension. The lower growth estimates
need no finite generation theorem: choose pluricanonical sections whose
ratios realize the Iitaka dimension in a fixed degree; independent
monomials in those ratios give the required polynomial lower bound in
multiples of that degree. $\square$

The hypothesis is concretely testable: $r_m$ global sections of
$E_m(b_mD)$ forming a generic basis yield \eqref{eq:114:injection}.
It is stronger than what the conjecture requests and has not been shown
to hold in general. In particular, when $K_Y$ is torsion one may have
$D=0$, leaving no ample twisting room at all.

\paragraph{Attempt: what sublinear ample twisting actually proves.}
A more plausible output of generic generation is an injection
\[
 \mathcal O_Y(-b_mH)^{\oplus r_m}\hookrightarrow E_m,
 \qquad b_m=o(m),
\]
for a fixed ample effective divisor $H$. This suffices when $K_Y$ is big.
Indeed, choose $q$ and $G\geq0$ with $qK_Y\sim H+G$. For $m=q\ell$,
\[
 mK_Y-b_mH\sim(\ell-b_m)H+\ell G.
\]
For sufficiently large $m$, $\ell-b_m\geq\ell/2$, so multiplication by
the section of $\ell G$ and ample section growth give
\[
 h^0(Y,mK_Y-b_mH)\geq c m^{\dim Y}.
\]
The previous proof then gives $\kappa(X)\geq\kappa(F)+\dim Y$.
This isolates, rather than re-proves from scratch, the geometric room
available in the established general-type-base case. \sref{114}{1}

That reasoning fails sharply for a non-big base. On an elliptic curve,
$K_Y\sim0$ and an effective divisor $H$ of positive degree satisfies
\[
 h^0(Y,mK_Y-H)=h^0(Y,-H)=0
\]
for \emph{every} $m$. Even constant pole cost $b_m=1=o(m)$ cannot be
absorbed. Thus ``the poles grow sublinearly'' is insufficient unless
their direction is also controlled. Lemma~1 is one possible replacement,
but its canonical-direction hypothesis is exactly an unresolved part
of this route.

\paragraph{Attempt: finite \'etale descent without a false invariants count.}
The following restricted case can be handled completely. Suppose a finite
\'etale cover $Y'\to Y$ trivializes the fibration:
\[
 X'=X\times_Y Y'\simeq F\times Y'.
\]
Then $\kappa(X)=\kappa(F)+\kappa(Y)$.

Here are the descent details. For a finite Galois \'etale cover
$\pi:Z\to W$ with group $G$, \'etaleness identifies $K_Z=\pi^*K_W$ and
\[
 R(W,K_W)=R(Z,K_Z)^G.
\]
For any element $s$ of the latter graded canonical ring, the monic
polynomial $\prod_{g\in G}(T-gs)$ has invariant coefficients. Hence
$R(Z,K_Z)$ is integral over its invariant subring. Their fraction fields
have equal transcendence degree, so their Kodaira dimensions agree.
If there are no positive-degree sections downstairs, the norm of a
nonzero homogeneous section upstairs would produce one, so this also
handles $\kappa=-\infty$. Passing through a Galois closure handles a
non-Galois finite \'etale cover. None of this asserts that invariant
sections occupy a fixed fraction of every individual graded piece.

On a product, the complete pluricanonical section space is the tensor
product of the two factor spaces. Its rational map is the Segre product
of the two pluricanonical maps wherever defined. Choosing a common
sufficiently divisible degree gives
$\kappa(F\times Y')=\kappa(F)+\kappa(Y')$, with the usual $-\infty$
convention. Applying \'etale invariance to $X'\to X$ and $Y'\to Y$ proves
the restricted statement. This is a self-contained recovery of a familiar
isotrivial situation, not a proof for arbitrary fibrations.

\paragraph{Stress test: positivity is not a supply of untwisted sections.}
Let $L$ be a non-torsion degree-zero line bundle on an elliptic curve.
It is nef and numerically trivial, yet $H^0(Y,L^m)=0$ for every $m>0$.
A nonzero section would have an effective zero divisor of degree zero,
hence no zeros, and would trivialize $L^m$, a contradiction. Consequently
one cannot replace the generic-generation hypothesis in Lemma~1 by
numerical semipositivity, even for line bundles. This example is not
asserted to be a pluricanonical direct image of a counterexample fibration:
the geometric origin of $E_m$ is exactly the additional structure a
successful proof must exploit. The direct-image methods in
Cao--P\u{a}un use much more than arbitrary numerical positivity.
\eref{114}{1}

Likewise, a generically chosen fibre basis can always be cleared of poles
with some ample twist depending on $m$, but that observation supplies
neither a sublinear bound on the twist nor support in a canonical divisor.
Also, replacing a fibration by a finite \emph{ramified} cover changes the
canonical divisor by ramification; the \'etale descent proof does not
apply verbatim to such a cover.

\paragraph{Outcome.}
\textbf{Outcome: Conditional reduction.}

The attempt proves a canonical-direction pole-cost criterion and explains
exactly why the apparently weaker sublinear ample-pole criterion only
works directly over a big-canonical base. It also gives a complete
finite-\'etale-product special case. What is missing is geometric control
of the pluricanonical direct images strong enough to remove poles on a
general non-big base. Assuming all good minimal models or abundance would
use another unresolved conjectural package, not complete this argument.
No novelty claim is made for the special cases or standard descent facts.

% Keep the short local bibliography together in the standalone reading copy.
\Needspace{14\baselineskip}
% END RESEARCH ADDITION 114
\subsection{Sources}
\begin{itemize}\raggedright
\item[\textnormal{[S1]}]\hypertarget{source:114:1}{} C. Birkar, Iitaka conjecture C\_\{n,m\} in dimension six, arXiv:0806.4413 [math.AG], 2008.
\url{https://arxiv.org/abs/0806.4413}.
{\small\textit{Catalogue formulation source.}}
% BEGIN NEW REFERENCES 114
\item[\textnormal{[E1]}]\hypertarget{extra:114:1}{} Junyan Cao and Mihai P\u{a}un,
\emph{Kodaira dimension of algebraic fiber spaces over abelian varieties},
Invent. Math. 207 (2017), 345--387; arXiv:1504.01095v3 (2016),
Theorem 1.1 and Sections 2--4.
\url{https://arxiv.org/html/1504.01095v3}.
\item[\textnormal{[E2]}]\hypertarget{extra:114:2}{} Houari Benammar Ammar,
\emph{Note on Iitaka Conjecture $C_{n,m}$}, arXiv:2510.06412v2
(12 December 2025), Theorem 1.2. Recent manuscript result, not used in
any proof in this notebook.
\url{https://arxiv.org/html/2510.06412v2}.
% END NEW REFERENCES 114
\end{itemize}

\end{document}
